(1) 定常方程 \(\theta''+\delta e^\theta=0\)。对偶解有 \(\theta'(0)=0\)，由自治常微分方程初值唯一性，正解写成
\[
\theta_a(x)=2\ln\frac{\cosh a}{\cosh(ax)},\quad a>0.
\]
直接求导得 \(\theta_a''=-2a^2\operatorname{sech}^2(ax)\)，且
\(e^{\theta_a}=\cosh^2 a\,\operatorname{sech}^2(ax)\)，因此
\[
\delta(a)=2a^2\operatorname{sech}^2a.
\]
该解在 \(x=\pm1\) 均为0，内部严格为正。任意给定正中心值的解由初值唯一性落在此族之内。

(2) 对函数求导：
\[
\delta'(a)=4a\operatorname{sech}^2a(1-a\tanh a).
\]
由于 \(a\tanh a\) 严格递增，唯一临界 \(a_c\tanh a_c=1\)；热爆炸阈值
\(\delta_c=2a_c^2\operatorname{sech}^2a_c\)。
\(0<\delta<\delta_c\) 有低温/高温两支，等号只有一支，超过则无正稳态。小 \(a\) 时 \(\delta=2a^2-2a^4+O(a^6)\)，大 \(a\) 时 \(\delta\sim8a^2e^{-2a}\)。过强的指数放热不能在固定边界温度下取得定常平衡。

(3) 线性扰动 \(v(x)e^{\lambda t}\) 满足
\[
\lambda v=v''+\delta(a)e^{\theta_a}v,\qquad v(\pm1)=0.
\]
设 \(L=-\partial_{xx}-\delta e^{\theta_a}\)；取
\(w=\partial_a\theta_a=2[\tanh a-x\tanh(ax)]\)，它在内部严格正且两端为0。对定常式微分可得
\(Lw=\delta'(a)e^{\theta_a}\)。
在临界 \(a_c\) 有严格正的零本征函数，故首特征值为0；低温分支 \(\delta'(a)>0\)，由正超解的Sturm比较原理首特征值严格为正，故增长率均负。高温分支 \(\delta'<0\)，以 \(w\) 为Rayleigh试验函数给 \(\langle w,Lw\rangle<0\)，故至少有一个 \(\lambda>0\) 的失稳模。

(4) 将方程在 \([-1,1]\) 上积分，
\[
\frac d{dt}\int_{-1}^1\theta dx
=\theta_x(1)-\theta_x(-1)+\delta\int_{-1}^1e^\theta dx.
\]
定常时壁面净导出热量 \(4a\tanh a\)；源项积分为
\(\delta\cosh^2a\int_{-1}^1\operatorname{sech}^2(ax)dx
=2a^2(2\tanh a/a)=4a\tanh a\)，严格匹配。小 \(a\) 时
\(\theta_a(x)=a^2(1-x^2)+O(a^4)\)，于是中心温升 \(\delta/2+O(\delta^2)\)，单侧壁面温度梯度模为 \(\delta+O(\delta^2)\)，吻合线性热传导近似。
